% translate_docs: manuals/preconditioning_manual.tex @ 9cb963f091f961327ab70fed56dbb0a05c4c6f5f
\documentclass[superscriptaddress,onecolumn,floatfix]{revtex4-2}

\usepackage{graphicx}
\usepackage{epsfig}
\usepackage{color}
\usepackage{amssymb}
\usepackage{amsmath}
\usepackage{array}
\usepackage[loose]{units}
\usepackage[russian]{babel}
\usepackage[unicode]{hyperref}
\usepackage{cleveref}
\usepackage{tasks}
\usepackage[caption=false]{subfig}
\usepackage{letltxmacro}

\LetLtxMacro{\ORIGselectlanguage}{\selectlanguage}
\makeatletter
\DeclareRobustCommand{\selectlanguage}[1]{%
    \@ifundefined{alias@\string#1}
      {\ORIGselectlanguage{#1}}
      {\begingroup\edef\x{\endgroup
         \noexpand\ORIGselectlanguage{\@nameuse{alias@#1}}}\x}%
}
\newcommand{\definelanguagealias}[2]{%
  \@namedef{alias@#1}{#2}%
}
\makeatother

\definelanguagealias{ru}{russian}


\begin{document}
\selectlanguage{ru}

\begin{abstract}
  В данном руководстве рассматриваются основные аспекты использования комплекта для ручной предварительной подготовки и часто задаваемые вопросы.
\end{abstract}
\title{Руководство к комплекту для предварительной подготовки}
\author{Бутик Electroniq Buttons}
\email{support@electroniqbuttons.com}
\date{\today}

\maketitle


\section{Какие автомобили могут это использовать?} \label{sect:eligibility}
Этот комплект предназначен для большинства автомобилей Ioniq 5 (2021–2024 гг.), Ioniq 6 (2022–2025 гг.) и EV6 (2022–2024 гг.). Мы провели испытания на автомобилях Ioniq 5 с трех разных рынков, а также на одном Ioniq 6 и одном EV6. \href{mailto:requests@electroniqbuttons.com}{Свяжитесь с нами} ...если вы планируете использовать этот комплект на другом автомобиле. Комплект будет работать только при наличии системы подогрева аккумуляторной батареи. Чтобы проверить это, поищите пункт меню «Battery Conditioning Mode» (Режим подготовки батареи) в разделе «EV». $\to$ «Настройки EV» (см. рис. \ref{subfig:menus:EVsettings}) на головном устройстве вашего автомобиля (центральном дисплее) (см. рис. \ref{subfig:menus:batt_cond}).
\begin{figure}[h]
  \subfloat[]{\label{subfig:menus:EVsettings}
  \includegraphics[height=1.95in]{photos/EV_settings.jpg} }
  \subfloat[]{\label{subfig:menus:batt_cond}
  \includegraphics[height=1.95in]{photos/battery_conditioning_mode.jpg} }
   \caption{(a) Как найти меню настроек электромобиля (EV): на главном экране нажмите значок «EV», затем в появившемся меню слева выберите пункт «Настройки EV» (EV Settings). (b) Экран режима подготовки аккумулятора (Battery Conditioning Mode). \label{fig:menus}}
\end{figure}

Если у вас нет такого экрана, этот комплект приобретать не следует. Пожалуйста, обратитесь к официальному дилеру или к опытным владельцам автомобилей в вашей стране, чтобы выяснить, можно ли установить эту функцию на ваш автомобиль с помощью дополнительных действий. Если вы находитесь в США, Канаде или Норвегии, вы можете связаться с нами по адресу \href{mailto:support@electroniqbuttons.com}{support@electroniqbuttons.com} и мы с радостью поможем вам оценить это, если вы предоставите немного информации о своем автомобиле.

Если вместо режима подготовки аккумулятора (Battery Conditioning Mode), показанного на рисунке \ref{subfig:menus:batt_cond}...у вас есть «Зимний режим» (Winter Mode), поэтому автомобилю требуется лишь обновление прошивки у официального дилера. Найдите соответствующий технический бюллетень для вашей страны, касающийся этого обновления, и обратитесь к ближайшему компетентному дилеру. После выполнения этой процедуры и подтверждения того, что пункт меню «Зимний режим» был заменен на «Режим подготовки аккумулятора» (Battery Conditioning Mode), вы сможете приобрести этот комплект.

\section{Базовая настройка комплекта}
Начните с \href{https://github.com/tylerharvey/Ioniq5_CAN/blob/main/guides/}{руководства по установке}. См. \href{https://meatpihq.github.io/wican-fw/config/wifi}{Документация WiCAN} для базовой настройки WiCAN. Это не требуется для активации функции предварительной подготовки (preconditioning), однако настоятельно рекомендуется как минимум изменить пароль для Wi-Fi, установленный по умолчанию, в целях безопасности.  

\section{Пользовательский интерфейс}
Мы активно работаем над пользовательским интерфейсом, поэтому заглядывайте сюда почаще!

Чтобы настроить параметры ручной предварительной подготовки, подключитесь к сети Wi-Fi устройства WiCAN и перейдите на страницу его настройки по адресу \href{http://192.168.80.1}{http://192.168.80.1}. Третья вкладка под названием «Buttons», показанная на рисунке \ref{subfig:ui:buttons} позволяет настроить кнопку ручной предварительной подготовки и любые другие пользовательские кнопки, добавленные впоследствии.

\begin{figure}[h]
  \subfloat[]{\label{subfig:ui:buttons}
  \includegraphics[height=1.95in]{photos/ui_2026-07-04.png} }  \\
   \caption{Снимок экрана страницы «Buttons» в HTTP-конфигурации WiCAN.}
\end{figure}

Первый вариант — «Режим предварительной подготовки» (Preconditioning Mode) — в настоящее время находится в разработке и предусматривает лишь одну настройку: однократное выполнение. После завершения одного цикла предварительной подготовки процесс прекращается. Мы также работаем над режимами «Непрерывный» (Continuous) и «Постоянный» (Persistent), которые позволят системе управления батареей (BMS) автоматически повторно нагревать аккумулятор при падении температуры: либо до момента выключения автомобиля (режим «Непрерывный»), либо даже с сохранением активности после циклов включения и выключения автомобиля (режим «Постоянный»).

Вторая опция позволяет выбрать кнопку, на которую будет назначена функция предварительной подготовки (preconditioning). Больше вариантов будет доступно при использовании настроенного WiCAN в режиме «человек посередине» (man-in-the-middle). По умолчанию отображаются 5 вариантов:
\begin{enumerate}
  \item Отключено
  \item Кнопка «Звезда» (на руле) — по умолчанию
  \item Кнопка «Звезда» (панель управления климатом)
  \item Нажатие на регулятор (панель климат-контроля)
  \item Пресс для придания объема (Climate Bar)
\end{enumerate}
При выборе настройки «Отключено» (Disabled) ни одна кнопка не будет запускать предварительную подготовку автомобиля. Это может быть полезно при сдаче автомобиля в сервисный центр дилера, чтобы не сбивать с толку обслуживающий персонал. Варианты с пометкой «нажатие» (press) подразумевают именно нажатие кнопок, в отличие от их поворота или перемещения вверх/вниз. Если ни один из этих вариантов не подходит, можно установить флажок «Показать все варианты кнопок» (Show all button options), чтобы открыть доступ к дополнительным настройкам:
\begin{tasks}(2)
  \task Режим (рулевое колесо)
  \task Голосовое управление (рулевое колесо)
  \task Вызов (рулевое колесо)
  \task Нажатие кнопки громкости (рулевое колесо)
  \task Увеличение громкости (руль)
  \task Уменьшение громкости (руль)
  \task Переход к следующему треку (рулевое колесо)
  \task Переход к следующему пункту вниз (рулевое колесо)
  \task OK (рулевое колесо)
  \task Карта (шкала климатических показателей)
  \task Панель навигации (панель управления климатом)
  \task Медиа (Climate Bar)
  \task Tuner Up (панель управления климатом)
  \task Настройка вниз (панель климат-контроля)
\end{tasks}
Выберите кнопку, которая вам подходит! Надеемся вскоре предложить и физические кнопки.

Чтобы запустить предварительную подготовку вручную, просто нажмите выбранную вами кнопку активации. На приборной панели сразу же появится обратный отсчет (в метрах или ярдах), временно заменяющий собой индикатор расстояния до пункта назначения. Если включен круиз-контроль, он перекрывает основной индикатор расстояния до цели. Используйте кнопку переключения страниц на рулевом колесе, чтобы перейти к экрану приборной панели, на котором отображаются навигационные подсказки. Процесс предварительной подготовки начнется, когда отсчет дойдет до нуля. Если отсчет начинается заново, это означает, что произошел какой-то сбой и система предпринимает повторную попытку. 

При активации предварительной подготовки вы увидите стандартное всплывающее сообщение «Предварительная подготовка аккумулятора активирована для оптимальной зарядки постоянным током», показанное на рисунке. \ref{subfig:precon:popup}. Также над символом аккумулятора в левом нижнем углу приборной панели появится значок снежинки или красной спирали, как показано на рисунке \ref{subfig:precon:coil}.

\begin{figure}[h]
  \subfloat[]{\label{subfig:precon:popup}
  \includegraphics[height=1.95in]{photos/precon_active.jpg} }  \\
  \subfloat[]{\label{subfig:precon:coil}
  \includegraphics[height=1in]{photos/coil_on_battery.jpg} }
   \caption{(a) Снимок экрана с всплывающим сообщением об успешной активации функции предварительной подготовки аккумулятора. (b) Пример снимка экрана, на котором значок спирали наложен на символ аккумулятора в левом нижнем углу приборной панели. \label{fig:precon}}
\end{figure}

Если предварительная подготовка не запускается, убедитесь, что соблюдены встроенные условия для ее активации:
\begin{enumerate}
  \item Температура аккумулятора должна быть ниже $\unit[21]{C}$/$\unit[70]{F}$
  \item Уровень заряда должен быть выше $\unit[20]{\%}$
  \item Экран «Режим подготовки аккумуляторной батареи» находится в меню «Настройки EV» (см. раздел \ref{sect:eligibility})
\end{enumerate}
К сожалению, это ограничения, заложенные в программное обеспечение блока управления аккумулятором или в аппаратную часть автомобиля, которые мы не можем изменить с помощью этого комплекта.

\section{Энергопотребление}
Этот комплект работает от бортовой сети автомобиля с напряжением 12 В, а значит, может потреблять энергию от штатного 12-вольтового аккумулятора. В связи с этим стоит обсудить несколько моментов:
\begin{enumerate}
  \item Сколько электроэнергии оно потребляет?
  \item Убьет ли это 12-вольтовый аккумулятор?
  \item MITM против параллельной схемы
\end{enumerate}

\subsection{Сколько электроэнергии оно потребляет?}

Мы измерили энергопотребление нескольких микроконтроллеров в различных режимах работы при выключенном двигателе (напряжение 12,7 В). Тестирование WiCAN в режиме «обычного использования» проводилось без выполнения HTTP-запросов и без подключения по BLE (например, через приложение CarScanner). Испытания выполнялись с использованием переключаемого жгута CAN-шины; насколько нам известно, передача OBD-запросов через него в данный момент невозможна, поэтому микроконтроллеры с поддержкой стандартного протокола OBD (WiCAN, vLinker MC+), вероятно, потребляли меньше энергии, чем при активном опросе параметров через штатный разъем OBD. 

\subsubsection{Оригинальный WiCAN с версией прошивки 4.20\_eb0.1.1-beta19}
\begin{enumerate}
  \item{Нормальная эксплуатация:} 30 мА
  \item{Использование HTTP-страницы:} до 34 мА
  \item{загрузка:} до 36 мА
  \item{Монитор шины CAN:} до 38 мА
  \item{Сон:} $<$1 мА
\end{enumerate}

\subsubsection{WiCAN (кастомная версия первой ревизии) с прошивкой v4.20\_eb0.1.1-beta19}
\begin{enumerate}
  \item{Стандартное использование MITM:} 40 мА
  \item{Обычное параллельное использование:} 34 мА
  \item{загрузка:} до 48 мА
\end{enumerate}

\subsubsection{WiCAN Pro с версией ПО 4.50}
\begin{enumerate}
  \item{Нормальная эксплуатация:} 50 мА
\end{enumerate}

\subsubsection{Белая панда бежит — последний аниматроник\_сборка «Панда»}
\begin{enumerate}
  \item{Нормальная эксплуатация:} 32 мА
\end{enumerate}

\subsubsection{vLinker MC+}
\begin{enumerate}
  \item{Обычное использование (без подключения):} 15–18 мА
  \item{CarScanner подключен через BLE:} 20–25 мА, всплески до 60 мА
\end{enumerate}

\subsection{Убьет ли это 12-вольтовый аккумулятор?}

В устройстве WiCAN предусмотрен режим энергосбережения (так называемый спящий режим), при котором отключается практически всё, а потребление энергии становится минимальным. Обычно этот режим настраивается на автоматическую активацию в зависимости от напряжения бортовой сети автомобиля (12 В). Если вы беспокоитесь о состоянии аккумулятора, рекомендуем включить эту функцию; найти её можно в интерфейсе WiCAN по адресу \href{http://192.168.80.1}{http://192.168.80.1} на вкладке «Энергосбережение». Дополнительную документацию можно найти здесь: \href{https://meatpihq.github.io/wican-fw/config/sleep-mode}{https://meatpihq.github.io/wican-fw/config/sleep-mode}Мы активно работаем над тем, чтобы перевести эту функцию на работу по шине CAN (то есть с отслеживанием сигнала выключения двигателя) вместо использования контроля напряжения. 

За 7 дней работы от нового 12-вольтового аккумулятора емкостью 60 А·ч типичное энергопотребление модифицированного устройства WiCAN (40 мА) привело бы к расходу $40 \unit[\mathrm{mA}] * 24 \unit[\mathrm{hours/day}] * 7 \unit[\mathrm{days}] = 6.7 Ah$, что составляет около 10\% of the 12V state of charge. Unlocked and off, a 2022 Canadian Preferred AWD Ioniq 5 with no aftermarket loads draws about 1.75 A (1750 mA). After ECUs quiet down 5 minutes after locking the car, the draw is still 0.75 A (750 mA). (The car has a further "sleep" mode after 7 days of inactivity; this draw hasn't been measured here.) So, a customized WiCAN in MITM mode and not in sleep mode increases draw over the first 7 days of inactivity by 5\%.

\subsection{MITM против параллельной схемы}
Если жгут проводов комплекта предпускового подогрева работает в режиме MITM (посредника), а устройство WiCAN переходит в спящий режим, запросы из приложений Hyundai/Kia не будут проходить. В наших планах на будущее — добавить функцию частичного сна, которая позволит передавать данные по шине CAN, но при этом отключит Wi-Fi. Если жгут используется в режиме MITM и вам необходимо отправлять запросы через приложения Hyundai/Kia, мы рекомендуем настроить WiCAN так, чтобы переход в спящий режим происходил при очень низком напряжении (например, 12,3 В). 

\section{Часто задаваемые вопросы}
\begin{enumerate}
  \item[В:] Как мне это установить?
  \item[A:] Посмотрите на \href{https://github.com/tylerharvey/Ioniq5_CAN/blob/main/guides/}{руководства по установке}Если это кажется слишком сложным, обратитесь за помощью к знакомому подростку или в мастерскую по установке автозвука!
  \item[В:] Почему камера заднего вида, автомобильное приложение или другая функция центрального дисплея не работают после установки?
  \item[A:] Проверьте \href{https://github.com/tylerharvey/Ioniq5_CAN/blob/main/guides/harnesses/head_unit_MITM/MITM_harness_modes.pdf}{Руководство по режимам работы жгута}Если у вас оригинальный WiCAN или в порт OBD не подключено никакое устройство, жгут должен находиться в параллельном режиме. Если у вас модифицированный WiCAN, жгут должен быть переведен в режим MITM.
  \item[В:] Станет ли из-за этого проще угнать мой автомобиль?
  \item[A:] Это маловероятно, однако устройство спроектировано так, чтобы его было трудно заметить — в том числе из соображений безопасности. Настоятельно рекомендуем сменить пароль для WiCAN. 
  \item[В:] Аннулирует ли это мою гарантию?
  \item[A:] Ознакомьтесь с законодательством вашей страны, регулирующим гарантийные обязательства и защиту прав потребителей. В США, например, закон Магнусона-Мосса (Magnuson-Moss Warranty Act) обеспечивает надежную защиту прав потребителей; угрозы автопроизводителей аннулировать гарантию из-за использования неоригинальных компонентов зачастую не выдерживают проверки в суде. (Данная информация не является юридической консультацией.) В то же время диагностика неисправностей специалистом автосервиса может быть затруднена, поэтому мы настоятельно рекомендуем демонтировать любые неоригинальные детали и аксессуары, способные осложнить процесс проверки. Например, при возникновении проблем с головным устройством или USB-портом перед визитом к официальному дилеру стоит снять данный комплект, чтобы избежать сложностей при диагностике. 
  \item[В:] Какова рекомендуемая комплектация набора для сдачи автомобиля в дилерский центр на базовое техническое обслуживание?
  \item[A:] Мы рекомендуем отсоединить WiCAN и убедиться, что жгут проводов установлен \href{https://github.com/tylerharvey/Ioniq5_CAN/blob/main/guides/harnesses/head_unit_MITM/MITM_harness_modes.pdf}{параллельный режим}.
  \item[В:] А что, если я захочу собрать это самостоятельно?
  \item[A:] Сделай это! Всё, что мы сделали — это \href{https://github.com/tylerharvey/Ioniq5_CAN}{с открытым исходным кодом}. Вы можете свободно использовать схему жгута проводов для самостоятельной сборки. Разъемы доступны на AliExpress. Однако, если вы внесете вклад в общедоступные знания в виде пулл-реквеста в \href{https://github.com/dragz/egmpdbc}{Файл DBC} ...или что-то иное — существуют программы возмещения части затрат, которые снижают итоговую стоимость до уровня цен на самостоятельное выполнение работ (DIY), при этом требуя гораздо меньше трудозатрат.
  \item[В:] Почему это нельзя сделать с помощью обновления программного обеспечения?
  \item[A:] Возможно. У нас есть свои догадки, почему Hyundai/Kia не включили это в обновление навигации. Если вы знаете, как... \href{https://discord.gg/zyrngQ5qQ}{взломать головное устройство} и нашли функции интерфейса CAN, вы можете их добавить \href{https://www.reddit.com/r/KiaEV6/comments/1jdd56u/comment/mid9xie/}{кнопка ручной активации предклиматизации для вашего головного устройства}Однако в настоящее время распространение взломанной прошивки невозможно, поэтому обеспечить широкую доступность такого неофициального «обновления» было бы затруднительно.
  \item[В:] Могу ли я реализовать это на своем устройстве Comma?
  \item[A:] Мы провели проверку и считаем, что это не работает. Похоже, сообщения о предварительной подготовке (preconditioning) не передаются с шин CAN A или E, к которым подключено устройство Comma, на шины M или G, к которым подключен блок BMU.
  \item[В:] Убьет ли это мой ICCU?
  \item[A:] Нет.
  \item[В:] Можно ли использовать встроенную навигацию при использовании этого комплекта?
  \item[A:] В бета-версии это невозможно. Это будет исправлено в модифицированной версии WiCAN.
  \item[В:] Повлияет ли обновление навигации на возможность использования комплекта?
  \item[A:] Это крайне маловероятно, поскольку работа комплекта не зависит от связи с головным устройством. Обновление блока управления аккумулятором могло бы повлиять на ситуацию (например, обновление, задним числом добавившее функцию предварительного прогрева батареи на ранних экземплярах Ioniq 5), но мы считаем, что такие обновления должны выполняться дилерами. Тем не менее, мы не можем гарантировать это со стопроцентной уверенностью.\% certainty that a navigation update will not affect the functionality of the kit. In the off chance that this does happen, let us know and we will do our best to update logic to restore manual preconditioning functionality.
  \item[В:] К чему этот обратный отсчет? Почему нельзя просто запустить предварительную подготовку сразу после нажатия кнопки?
  \item[A:] Блок управления аккумулятором (ЭБУ внутри батареи, содержащий BMS) имеет встроенную задержку в 60 секунд после получения первоначального запроса на предварительную подготовку. Если запрос прекращается до истечения этого 60-секундного интервала, процесс предварительной подготовки не запускается. Поэтому мы постарались сделать так, чтобы пользователю было очевидно и понятно, что необходимо подождать эти 60 секунд.
  \item[В:] Вы сделаете для меня новую кнопку?
  \item[A:] Возможно! Напишите нам на \href{mailto:requests@electroniqbuttons.com}{requests@electronicbuttons.com} с вашими запросами.
  \item[В:] Ваш бизнес называется «Electroniq Buttons Boutique». А где же физические кнопки?
  \item[A:] Мы работаем над тем, чтобы сделать доступными физические кнопки. У Роя есть протестированный прототип, показанный в его \href{https://github.com/dragz/ironiq}{Ironiq} репозиторий.
\end{enumerate}
% \begin{figure}[h]
%   \subfloat[]{\label{subfig:overview:harness}
%   \includegraphics[height=1.95in]{photos/whole_harness.jpg} }
%   \subfloat[]{\label{subfig:overview:shunt_cap_labeled}
%   \includegraphics[height=1.95in]{photos/labeled_connectors_MITM_harness.jpg} }
%    \caption{(a) Man-in-the-middle harness. (b) Overview of connectors showing shunt cap. \label{fig:overview}}
% \end{figure}
\appendix 
\bibliographystyle{apsrev4-1}
\bibliography{stemh_model}{}
\end{document}